#include "main.typ"
\section{Kết luận}
\addcontentsline{toc}{section}{Kết luận}

Chuyên đề \textbf{Cực trị của hàm số trùng phương} $y = ax^4 + bx^2 + c \ (a \neq 0)$ là một trong những mảng kiến thức then chốt, xuất hiện với tần suất cao trong các đề thi Tốt nghiệp THPT và thi đánh giá năng lực. Việc làm chủ kiến thức chuyên đề này không chỉ đòi hỏi sự thành thạo về mặt đại số mà còn yêu cầu khả năng tư duy hình học linh hoạt khi xử lý các bài toán liên quan đến tam giác cực trị $\Delta ABC$.

\vspace{0.5em}
\noindent \textbf{1. Ý nghĩa của việc ứng dụng công thức giải nhanh}
\begin{itemize}
    \item \textbf{Tối ưu hóa thời gian:} Việc áp dụng chính xác các công thức giải nhanh giúp giảm đáng kể số lượng bước biến đổi đại số cồng kềnh, từ đó tiết kiệm tối đa thời gian làm bài trắc nghiệm.
    \item \textbf{Hạn chế sai sót:} Các bài toán cực trị chứa tham số $m$ thường dẫn đến các hệ phương trình bậc cao dễ gây nhầm lẫn dấu hoặc thiếu điều kiện. Công thức tổng quát giúp định hướng lời giải chính xác và gãy gọn.
\end{itemize}

\vspace{0.5em}
\noindent \textbf{2. Tóm tắt các nguyên tắc cốt lõi cần tuân thủ}
\begin{itemize}
    \item \textbf{Kiểm tra điều kiện $a \neq 0$:} Luôn luôn xét trường hợp $a = 0$ trước khi áp dụng các công thức của hàm bậc bốn trùng phương nếu hệ số $a$ có chứa tham số.
    \item \textbf{Điều kiện tồn tại 3 cực trị ($ab < 0$):} Đây là tiền đề bắt buộc trước khi triển khai các bài toán về góc, diện tích, bán kính đường tròn hay độ dài cạnh của tam giác cực trị.
    \item \textbf{Khai thác tính chất hình học:} Điểm cực trị $A(0; c)$ luôn nằm trên trục tung $Oy$ và tam giác $ABC$ luôn cân tại $A$. Tận dụng tính chất đối xứng này sẽ giúp giải quyết nhanh các bài toán liên quan đến tọa độ trọng tâm, tâm đường tròn ngoại tiếp hoặc khoảng cách.
\end{itemize}

\vspace{0.5em}
\noindent \textbf{3. Lời khuyên cho quá trình ôn luyện}
\begin{itemize}
    \item \textbf{Hiểu bản chất trước khi học thuộc:} Người học nên chủ động nắm vững cách chứng minh các công thức giải nhanh thay vì học vẹt, nhằm linh hoạt xử lý khi đề bài biến tấu dạng hỏi.
    \item \textbf{Kết hợp phương pháp tự luận và trắc nghiệm:} Thành thạo cả hai phương pháp giúp người học vừa xây dựng tư duy toán học vững chắc để làm bài tự luận, vừa đảm bảo tốc độ và độ chính xác tuyệt đối trong các kỳ thi trắc nghiệm.
\end{itemize}

\vspace{1em}
Hy vọng tài liệu tổng hợp này sẽ là công cụ hỗ trợ đắc lực, giúp các bạn học sinh củng cố toàn bộ lý thuyết, làm chủ các dạng toán và tự tin đạt kết quả cao nhất trong các kỳ thi sắp tới.
